\documentclass{beamer}

\input{settings.tex}


\title{H infinity, Young}
\subtitle{\mycoursetitle, Lecture 10}
\author{by Sergei Savin}
\centering
\date{\mydate}



\begin{document}
\maketitle



\begin{frame}{\texorpdfstring{$\mathcal{H}_\infty$}{} Control, Storage Functions, and KYP}
	\begin{flushleft}
		
		{\small
		\textbf{System:}
		\[
		\dot{x} = Ax + Bw, \quad y = Cx + Dw
		\]
		
		
		\textbf{$\mathcal{H}_\infty$ objective:} ensure
		\[
		\|y\|_2 \le \gamma \|w\|_2 \quad \text{for all } w \in \mathcal{L}_2
		\]
		
		
		\textbf{Storage function (dissipation inequality):} \\
		Find $V(x) = x^\top P x$, $P \succ 0$, such that
		\[
		\dot{V}(x) + y^\top y - \gamma^2 w^\top w \le 0
		\]
		
		\textbf{KYP (Bounded Real Lemma):} \\
		The above holds iff $\exists P \succ 0$ such that
		\[
		\begin{bmatrix}
			A^\top P + PA & PB & C^\top \\
			B^\top P & -\gamma^2 I & D^\top \\
			C & D & -I
		\end{bmatrix} \prec 0
		\]
		
		
		\textbf{Conclusion:} $\Rightarrow \|G(s)\|_\infty < \gamma$}
		
	\end{flushleft}
	
\end{frame}


\begin{frame}{H infinity}
	\begin{flushleft}
		
Consider the uncertain system:
\begin{align*}
	\dot{x} &= (A + \Delta A)x + Bu + Dw, \\
	y &= Cx,
\end{align*}
with structured uncertainty
\[
\Delta A = N \Delta M, \quad \|\Delta\| < 1.
\]

We apply state feedback:
\[
u = Kx,
\]
yielding the closed-loop system:
\[
\dot{x} = (A + BK + \Delta A)x + Dw.
\]
		
	\end{flushleft}
\end{frame}



\begin{frame}{Step 1: KYP Lemma}
	\begin{flushleft}
		
	The $\mathcal{H}_\infty$ condition is satisfied if there exists $P \succ 0$ such that:
	\[
	\begin{bmatrix}
		(A + BK + \Delta A)^\top P + P(A + BK + \Delta A) & PD & C^\top \\
		D^\top P & -\gamma I & 0 \\
		C & 0 & -\gamma I
	\end{bmatrix} \prec 0.
	\]
	
		
	\end{flushleft}
\end{frame}




\begin{frame}{Step 2: Separate nominal and uncertain parts}
	\begin{flushleft}
		
		
		Let:
		\[
		W_1 =
		\begin{bmatrix}
			(A + BK)^\top P + P(A + BK) & PD & C^\top \\
			D^\top P & -\gamma I & 0 \\
			C & 0 & -\gamma I
		\end{bmatrix},
		\]
		\[
		W_2 =
		\begin{bmatrix}
			(\Delta A)^\top P + P(\Delta A) & 0 & 0 \\
			0 & 0 & 0 \\
			0 & 0 & 0
		\end{bmatrix}.
		\]
		
		Using $\Delta A = N \Delta M$, we obtain:
		\[
		W_2 =
		\begin{bmatrix}
			M^\top \Delta^\top N^\top P + PN \Delta M & 0 & 0 \\
			0 & 0 & 0 \\
			0 & 0 & 0
		\end{bmatrix}.
		\]
		
	\end{flushleft}
\end{frame}


\begin{frame}{Young's Inequality (Matrix Form)}
	
	\begin{flushleft}
		
		\textbf{Statement:} For any matrices $X, Y$ of compatible dimensions and any $\varepsilon > 0$,
		\[
		X^\top Y + Y^\top X \;\preceq\; \varepsilon X^\top X + \varepsilon^{-1} Y^\top Y
		\]
		
		
		\medskip
		
		\textbf{Equivalent form (used in robust control):}
		\[
		X \Delta Y + (X \Delta Y)^\top
		\;\preceq\;
		\varepsilon XX^\top + \varepsilon^{-1} Y^\top Y,
		\quad \|\Delta\| \le 1
		\]
		
		\medskip
		
		\textbf{Role in LMIs:}
		\begin{itemize}
			\item Replaces uncertain terms with \emph{convex bounds}
			\item Introduces scalar multiplier $\varepsilon$
		\end{itemize}
		
	\end{flushleft}
	
\end{frame}



\begin{frame}{Factorization of Uncertain Term}
	\begin{flushleft}
		
		
		
		\textbf{Uncertain block:}
		\[
		W_2 =
		\begin{bmatrix}
			M^\top \Delta^\top N^\top P + PN \Delta M & 0 & 0 \\
			0 & 0 & 0 \\
			0 & 0 & 0
		\end{bmatrix}
		\]
		
		\medskip
		
		\textbf{Define:}
		\[
		V_1 =
		\begin{bmatrix}
			PN \\ 0 \\ 0
		\end{bmatrix},
		\quad
		V_2 =
		\begin{bmatrix}
			M^\top \\ 0 \\ 0
		\end{bmatrix}
		\]
		
		\medskip
		
		\textbf{Compact factorization:}
		\[
		W_2 = V_1 \Delta V_2^\top + V_2 \Delta^\top V_1^\top
		\]
		
		\medskip
		
		\textbf{Expanded form:}
		\[
		W_2 =
		\begin{bmatrix}
			PN \\ 0 \\ 0
		\end{bmatrix}
		\Delta
		\begin{bmatrix}
			M & 0 & 0
		\end{bmatrix}
		+
		\begin{bmatrix}
			M^\top \\ 0 \\ 0
		\end{bmatrix}
		\Delta^\top
		\begin{bmatrix}
			N^\top P & 0 & 0
		\end{bmatrix}
		\]
	
	\end{flushleft}
	
\end{frame}



\begin{frame}{Step 3: Young's inequality}
	\begin{flushleft}

		
		Using Young's inequality:
		\[
		W_2 = V_1 \Delta V_2^\top + V_2 \Delta^\top V_1^\top \preceq \frac{1}{\varepsilon} V_1 V_1^\top + \varepsilon V_2 V_2^\top,
		\]
		we replace the condition by:
		\[
		W_1 + \frac{1}{\varepsilon} V_1 V_1^\top + \varepsilon V_2 V_2^\top \prec 0.
		\]
		
		
	\end{flushleft}
\end{frame}


\begin{frame}{Step 4: First Schur complement}
\begin{flushleft}
	
	Applying Schur complement to $\frac{1}{\varepsilon} V_1 V_1^\top$:
	\[
	\begin{bmatrix}
		W_1 + \varepsilon V_2 V_2^\top & V_1 \\
		V_1^\top & -\varepsilon I
	\end{bmatrix}
	\prec 0.
	\]
	
	Expanding:
	\[
	\begin{bmatrix}
		(A + BK)^\top P + P(A + BK) + \varepsilon M^\top M & PD & C^\top & PN \\
		D^\top P & -\gamma I & 0 & 0 \\
		C & 0 & -\gamma I & 0 \\
		N^\top P & 0 & 0 & -\varepsilon I
	\end{bmatrix}
	\prec 0.
	\]
	
\end{flushleft}
\end{frame}


\begin{frame}{Step 5: Change of variables}
\begin{flushleft}
	
	Let:
	\[
	Q = P^{-1}, \quad U = KQ.
	\]
	
	Apply the congruence transformation $\mathrm{diag}(Q, I, I, I)$:
	\[
	\begin{bmatrix}
		QA^\top + AQ + BU + U^\top B^\top + \varepsilon Q M^\top M Q
		& D & QC^\top & N \\
		D^\top & -\gamma I & 0 & 0 \\
		CQ & 0 & -\gamma I & 0 \\
		N^\top & 0 & 0 & -\varepsilon I
	\end{bmatrix}
	\prec 0.
	\]
	
	
\end{flushleft}
\end{frame}



\begin{frame}{Step 6: Second Schur complement}
	\begin{flushleft}
		
		\[
		\varepsilon Q M^\top M Q
		=
		(Q M^\top)(\varepsilon I)(M Q).
		\]
		
		We now apply the Schur complement:
		\[
		\begin{bmatrix}
			QA^\top + AQ + BU + U^\top B^\top
			& D
			& QC^\top
			& N
			& Q M^\top \\
			D^\top
			& -\gamma I
			& 0
			& 0
			& 0 \\
			CQ
			& 0
			& -\gamma I
			& 0
			& 0 \\
			N^\top
			& 0
			& 0
			& -\varepsilon I
			& 0 \\
			M Q
			& 0
			& 0
			& 0
			& -\varepsilon^{-1} I
		\end{bmatrix}
		\prec 0.
		\]
		
	\end{flushleft}
\end{frame}



\begin{frame}{Step 7: Change of variables}
	\begin{flushleft}
		
		Define:
		\[
		\alpha = \varepsilon^{-1}.
		\]
		
		Then multiply the LMI by the congruence transformation:
		\[
		\mathrm{diag}(I, I, I, \alpha I, I),
		\]
		to obtain:
		
		\[
		\begin{bmatrix}
			QA^\top + AQ + BU + U^\top B^\top
			& D
			& QC^\top
			& \alpha N
			& Q M^\top \\
			D^\top
			& -\gamma I
			& 0
			& 0
			& 0 \\
			CQ
			& 0
			& -\gamma I
			& 0
			& 0 \\
			\alpha N^\top
			& 0
			& 0
			& -\alpha I
			& 0 \\
			M Q
			& 0
			& 0
			& 0
			& -\alpha I
		\end{bmatrix}
		\prec 0.
		\]
		
	\end{flushleft}
\end{frame}



\begin{frame}{Conclusion}
	\begin{flushleft}
		
		The condition is linear in decision variables:
		\[
		Q \succ 0, \quad U, \quad \alpha  > 0, \quad \gamma>0,
		\]
		with controller recovered as:
		\[
		K = U Q^{-1}.
		\]
		
		
	\end{flushleft}
\end{frame}


\myqrframe


\end{document}
